\section{Algebraic cables of the trefoil}\label{sec:cable}
Beyond torus knots, refined $\SO(2N)$ invariants have had no definition, since the $S$ and $T$ matrices of refined Chern--Simons theory only produce torus knots. For iterated torus knots, the iterated DAHA-Jones polynomial of \cite{CD14} is defined for any reduced root system and weight $b$ by
\begin{equation}\label{eq:iterJD}
JD_{\vec r,\vec s}(b) = \bigl\{\hat\gamma_1\bigl(\cdots\hat\gamma_{\ell-1}\bigl((\hat\gamma_\ell(P_b)/P_b(q^{-\rho_k}))\!\Downarrow\bigr)\!\Downarrow\cdots\bigr)\bigr\}_{\mathrm{ev}},\qquad H\!\Downarrow\, := H\cdot 1,
\end{equation}
where $\gamma_\ell$ is the innermost pattern, $\gamma_1$ the companion, $\gamma_i$ lifts an element of $\mathrm{SL}(2,\mathbb{Z})$ with first column $(r_i, s_i)$, and $\{\cdot\}_{\mathrm{ev}}$ is the evaluation at $q^{-\rho_k}$; by Theorem 2.4 of \cite{CD14} the result does not depend on the choice of lifts. We adopt its $D_N$ value as the definition of the refined $\SO(2N)$ invariant of an algebraic cable.
\begin{hypothesis}\label{hyp:cablecover}
The cover of an algebraic cable is given by the $\GL_N$ iterated DAHA-Jones polynomials in ordinary and composite representations, normalized by \eqref{eq:Mb} (Hypothesis \ref{hyp:cover} for cables).
\end{hypothesis}
The examples are the $(2, q_2)$-cables of the trefoil, $\Cab{13}$ ($q_2 = 13$) and $\Cab{15}$ ($q_2 = 15$), with $q_2$ measured in the Seifert framing of the companion. The pattern word has first column $(2, q_2 - 12)$, i.e., $\tau_+\tau_-$ for $(2,1)$ and $\tau_-\tau_+\tau_-$ for $(2,3)$ in the notation of Appendix \ref{app:daha-conv}, and the companion is the trefoil $\tau_+\tau_-^2$ (the cable is algebraic if and only if $q_2 > 12$). Non-algebraic cables are not treated.
In the companion step the polynomial $f_1 = \hat\gamma_{\mathrm{pat}}(P_b)\cdot 1$ from the pattern step is applied as $f_1(\hat\gamma_{\mathrm{comp}}(X))\cdot 1$. For $D_N$, $f_1$ is invariant under the Weyl group of $D_N$ and, for the degrees relevant here ($\deg f_1 \le 2|b| < N$), also under that of $B_N$; it is expanded in a basis $H_\nu$ of $B_N$-symmetric Laurent polynomials indexed by partitions $\nu$ with $|\nu| \le \deg f_1$, so that $JD$ becomes a linear combination of the basis values $\{\hat\gamma_{\mathrm{comp}}(H_\nu)\}$, computed once per $N$ and shared between the two cables (the companion is common). The adjoint of the cover is expanded likewise in $H_\mu(X)H_\nu(X^{-1})$.

% [D15]
\begin{result}[normalization and general $N$ for cables]\label{cr:cable-interp}
For the cables $\Cab{13}$ and $\Cab{15}$ the following hold.
\begin{enumerate}
\item The $\GL_N$ values agree up to a monomial with the A-type values of the same cables obtained earlier by the author with an independent code ($\U(2)$, $\U(3)$, $|b| \le 2$).
\item The $D_N$ values differ from the mirror of the fixed-$N$ $W_R$ in the uniform normalization, built through $\SO(4) = \mathrm{SU}(2)^2$ and $\SO(6) = \mathrm{SU}(4)$, only by $T_R^{-F}$ with $F = p_2q_2$ ($26$ for $\Cab{13}$, $30$ for $\Cab{15}$), for $N = 2, 3$ in the fundamental representation. This is the torus-knot rule with $F = rs$ replaced by $p_2q_2$. The genus $\delta$ of \eqref{eq:uniform} is $8$ ($\Cab{13}$) and $9$ ($\Cab{15}$), from the genus formula $\delta \mapsto p_2\delta + (p_2-1)(q_2-1)/2$ of a $(p_2,q_2)$-cable.
\item On the $\GL_N$ side, \eqref{eq:Mb} with $F = p_2q_2$ and $\cK = 10$ ($\Cab{13}$), $12$ ($\Cab{15}$) holds for six ordinary representations at $N = 2$--4; the constants $\cK$ are empirical. Since the adjoint has $|b| = 0$, this fixes its normalization to $T_{\mathrm{adj}}^{-F}$.
\item The interpolation in $N$ (Appendix \ref{app:interp}, Table \ref{tab:cableinterp}) predicted the values at the $N$ not used coefficient by coefficient. For $(1,1)$ of $\Cab{13}$ the cancellations at small $N$ are strong; with the data for $N \le 13$ three degrees in $q^{1/2}$ remained undetermined, and adding $N = 14$ determined them. The check of this solution is at $N = 2, 3$ (all coefficients), which are the points where the cancellations are strongest, so it is weaker than for the other rows. The $(1,1)$ of $\Cab{15}$ requires $N \ge 13$ and the largest basis element $(1,1,1,1)$ (estimated above 60\,GB of memory), so the two-box computation was done for $\Cab{13}$ only.
\end{enumerate}
\end{result}

% [D16]
\begin{result}[crosscap part of cables at large $N$]\label{cr:cable-largeN}
The crosscap parts of the cables have the following properties.
\begin{enumerate}
\item Fundamental representation ($\Cab{13}$, $\Cab{15}$): $g_{\yone}$ is an integer Laurent polynomial and its refined part is divisible by the universal factor $(q-t)(\lambda^2-t)$, with 6 of 9 ($\Cab{13}$) and 8 of 11 ($\Cab{15}$) $\ba$ blocks of the quotient aligned. At $t = q$ the $c = 1$ part is symmetric under $q \leftrightarrow 1/q$, and the $c = 2$ part does not vanish (the unrefined $c = 2$ invariants of a cable are non-zero).
\item Two boxes ($\Cab{13}$): the structure is that of Result \ref{cr:largeN}(1), and the refined parts of $\hat g_{\yone}, \hat g_{\ytwo}, \hat g_{\yoneone}$ and of their sum are all divisible by the universal factor. The quotient for the two-box sum has 30 of 50 $\ba$ blocks aligned (Table \ref{tab:signs}).
\item The sign property (Definition \ref{def:signprop}) fails. ($c{=}1$): the $c = 1$ part is not $\sigma$-symmetric; only the extreme blocks ($\lambda^{-26}, \lambda^{-16}$ for $\Cab{13}$; $\lambda^{-30}, \lambda^{-28}, \lambda^{-18}$ for $\Cab{15}$) are symmetric, with aligned signs in their $G$-expansions. Splitting the $\sigma$-asymmetric part by Lefschetz weight, each block has a common sign in all rows, alternating with $\beta$, as for $T(3,4)$. ($c{=}2$): the refined part of the $c = 2$ part is divisible by the universal factor, but the quotient has mixed signs within the $\ba$ blocks (3 of 4 blocks mixed for $\Cab{13}$, 3 of 5 for $\Cab{15}$); the whole $c = 2$ part is not divisible by $(\bt^2 - 1)$, as it must not be since its unrefined limit is non-zero. The refined part of the $c = 1$ part, by contrast, has aligned signs in every $\ba$ block (5 blocks for $\Cab{13}$, 6 for $\Cab{15}$); the counts 6/9 and 8/11 of item 1 are the sums over the two parts.
\end{enumerate}
\end{result}
Thus for the cables what holds is integrality and the universal factor (Conjecture \ref{conj:structure}), while the sign property fails in both of its parts, in contrast to $T(3,4)$ where only ($c{=}1$) fails (\S\ref{sec:t34}). Whether the half-degree relation of \S\ref{sec:sigmaorigin} between the $c = 1$ part and the brane--image sector is modified for these knots has not been examined \open. The normalization of the cover for the cables rests on Hypothesis \ref{hyp:cablecover} and on the empirical constants $\cK$; the integrality and the universal factor at general $N$, which were not used in fixing them, are the checks of this choice.
